\documentclass[12pt,a4paper]{report}

\usepackage[brazil]{babel}
\usepackage[utf8]{inputenc}
\usepackage[T1]{fontenc}
\usepackage[left=2.5cm,right=2.5cm,top=2.5cm,bottom=2.5cm]{geometry}
\usepackage{graphicx}
\usepackage{amsmath,amssymb}
\usepackage{booktabs,multirow,array,longtable}
\usepackage[hidelinks]{hyperref}
\usepackage{xcolor}
\usepackage{fancyhdr}

\pagestyle{fancy}
\fancyhf{}
\lhead{Relatório Técnico}
\rhead{\thepage}
\renewcommand{\headrulewidth}{0.4pt}

\begin{document}

% ---------- CAPA ----------
\begin{titlepage}
\centering
\vspace*{2cm}

{\Large \textbf{NOME DA INSTITUIÇÃO / UNIVERSIDADE}}\\[1.5cm]

{\LARGE \textbf{RELATÓRIO TÉCNICO}}\\[0.7cm]
{\large \textbf{Título do Relatório}}\\[2.0cm]

\begin{flushleft}
\textbf{Autor:} Nome Completo\\
\textbf{Curso:} Nome do Curso\\
\textbf{Disciplina / Projeto:} Nome da disciplina ou projeto\\
\textbf{Professor / Orientador:} Nome do professor\\
\end{flushleft}

\vfill
Brasília -- DF\\
2025
\end{titlepage}

\pagenumbering{roman}
\tableofcontents
\listoffigures
\listoftables

\clearpage
\pagenumbering{arabic}

% ---------- TEXTO ----------
\chapter{Introdução}
Este relatório técnico apresenta uma visão estruturada sobre o tema estudado, com foco em organização, clareza e padronização da documentação. O objetivo deste modelo é servir como base para relatórios acadêmicos, projetos de engenharia e registros técnicos.

\chapter{Objetivos}
\section{Objetivo geral}
Descrever o problema, o contexto e os resultados obtidos no desenvolvimento do trabalho.

\section{Objetivos específicos}
\begin{itemize}
    \item Organizar a documentação técnica de forma padronizada.
    \item Apresentar dados, análises e resultados com clareza.
    \item Utilizar recursos do \LaTeX{} para melhor estruturação do texto.
\end{itemize}

\chapter{Metodologia}
A metodologia adotada pode incluir levantamento bibliográfico, análise de documentos, desenvolvimento de protótipos, testes e consolidação dos resultados.

\section{Procedimentos}
\begin{enumerate}
    \item Definição do problema.
    \item Levantamento de referências.
    \item Execução das atividades.
    \item Consolidação e análise dos resultados.
\end{enumerate}

\section{Exemplo de figura}
\begin{figure}[h]
    \centering
    \fbox{\rule{0pt}{4cm}\rule{0.8\textwidth}{0pt}}
    \caption{Placeholder para inserção de figura}
    \label{fig:placeholder}
\end{figure}

\section{Exemplo de tabela}
\begin{table}[h]
\centering
\caption{Exemplo de tabela com células mescladas}
\label{tab:exemplo}
\begin{tabular}{lccc}
\toprule
\multirow{2}{*}{Categoria} & \multicolumn{3}{c}{Indicadores} \\
\cmidrule(lr){2-4}
 & 2023 & 2024 & 2025 \\
\midrule
Qualidade & 85 & 90 & 94 \\
Desempenho & 78 & 83 & 88 \\
Acurácia & 80 & 86 & 91 \\
\bottomrule
\end{tabular}
\end{table}

\chapter{Resultados e Discussão}
Os resultados devem ser descritos de maneira objetiva, destacando tendências, impactos e observações relevantes. A discussão deve conectar os dados obtidos ao problema proposto.

\chapter{Conclusão}
O relatório evidencia a importância de uma estrutura clara, reprodutível e tecnicamente consistente. O uso do \LaTeX{} facilita a organização do conteúdo e melhora a apresentação final do documento.

\begin{thebibliography}{9}
\bibitem{lamport}
LAMPORT, Leslie. \textit{LaTeX: A Document Preparation System}. Reading: Addison-Wesley, 1994.

\bibitem{knuth}
KNUTH, Donald E. \textit{The TeXbook}. Reading: Addison-Wesley, 1984.
\end{thebibliography}

\end{document}